% F6 · Árbol de decisión: destino de cada conjunto (D13·01)
\documentclass[border=10pt]{standalone}
% ------------------------------------------------------------------
% Preámbulo común de las figuras del curso C1
% Paleta y tipografía alineadas con ../../estilo.tex
% ------------------------------------------------------------------
\usepackage[utf8]{inputenc}
\usepackage[T1]{fontenc}
\usepackage[spanish,es-noquoting]{babel}
\usepackage{helvet}
\renewcommand{\familydefault}{\sfdefault}
\usepackage{xcolor}
\usepackage{tikz}
\usetikzlibrary{arrows.meta,positioning,decorations.pathreplacing,calc,shapes.geometric,fit,backgrounds}

\definecolor{cHead}{RGB}{68,84,106}
\definecolor{cAcc}{RGB}{0,102,153}
\definecolor{cGris}{gray}{0.42}
\definecolor{cSuave}{RGB}{240,244,248}
\definecolor{cBien}{RGB}{34,124,68}
\definecolor{cMal}{RGB}{176,42,42}
\definecolor{cBienSuave}{RGB}{230,243,234}
\definecolor{cMalSuave}{RGB}{250,235,233}
\definecolor{cAccSuave}{RGB}{225,238,246}

\tikzset{
  caja/.style     = {rectangle, rounded corners=2.5pt, draw=cHead, line width=.7pt,
                     fill=cSuave, align=center, inner sep=6pt, font=\small},
  cajaAcc/.style  = {caja, draw=cAcc, fill=cAccSuave},
  cajaOk/.style   = {caja, draw=cBien, fill=cBienSuave},
  cajaNo/.style   = {caja, draw=cMal, fill=cMalSuave, dashed},
  cajaTenue/.style= {caja, draw=cGris, fill=white},
  rotulo/.style   = {font=\scriptsize\itshape, text=cGris, align=center},
  etiq/.style     = {font=\scriptsize, text=cHead, align=center,
                     fill=white, inner sep=1.5pt},
  fl/.style       = {-{Stealth[length=4.5pt]}, draw=cHead, line width=.7pt},
  flAcc/.style    = {fl, draw=cAcc},
  flNo/.style     = {-{Stealth[length=4.5pt]}, draw=cMal, line width=.7pt, dashed},
  llave/.style    = {decorate, decoration={brace, amplitude=4pt, raise=2pt},
                     draw=cGris, line width=.6pt},
  titulo/.style   = {font=\small\bfseries, text=cHead},
}

% Cabecera de figura: #1 número, #2 título
\newcommand{\figtit}[2]{\node[font=\footnotesize\bfseries, text=cGris, anchor=west]}

\begin{document}
\begin{tikzpicture}

\tikzset{
  preg/.style = {caja, text width=4.8cm, minimum height=1.05cm, font=\scriptsize},
  dest/.style = {text width=3.5cm, minimum height=1cm, font=\scriptsize\bfseries,
                 rounded corners=2.5pt, align=center, inner sep=6pt, line width=.8pt},
  si/.style   = {etiq, text=cBien, font=\scriptsize\bfseries},
  no/.style   = {etiq, text=cMal,  font=\scriptsize\bfseries},
}

% ---------------- columna de preguntas ------------------------------
\node[cajaTenue, text width=3.6cm, font=\scriptsize\bfseries] (ini) at (0,0)
  {Un conjunto del inventario};
\node[preg] (p1) at (0,-2.0) {¿Contiene datos personales?};
\node[preg] (p2) at (0,-4.0) {¿Lo permiten el consentimiento\\ y las autorizaciones?};
\node[preg] (p3) at (0,-6.0) {¿Existe un procedimiento de acceso\\ con órgano, criterios y plazo?};

\draw[fl] (ini) -- (p1);
\draw[fl] (p1) -- (p2) node[no, pos=.5] {sí};
\draw[fl] (p2) -- (p3) node[si, pos=.5] {sí};

% ---------------- destinos ------------------------------------------
\node[dest, draw=cBien, fill=cBienSuave, text=cBien] (ab) at (7.6,-2.0) {ABIERTO};
\node[dest, draw=cMal,  fill=cMalSuave,  text=cMal]  (nc) at (7.6,-4.0) {NO SE COMPARTE};
\node[dest, draw=cAcc,  fill=cAccSuave,  text=cAcc]  (ac) at (7.6,-6.0) {ACCESO\\ CONTROLADO};

\draw[fl, draw=cBien] (p1) -- (ab) node[si, pos=.45] {no};
\draw[fl, draw=cMal]  (p2) -- (nc) node[no, pos=.45] {no};
\draw[fl, draw=cAcc]  (p3) -- (ac) node[si, pos=.45] {sí};

% la respuesta «no» de la tercera pregunta no cierra el conjunto: obliga a definirlo
\node[cajaAcc, text width=5.4cm, font=\scriptsize] (def) at (0,-8.1)
  {Hay que \textbf{definirlo} antes de depositar.\\ No tenerlo no es motivo para cerrar el conjunto};
\draw[fl, draw=cAcc] (p3) -- (def) node[no, pos=.5] {no};
\draw[fl, draw=cAcc] (def.east) -| (ac.south);

% ---------------- los tres frenos -----------------------------------
\node[cajaTenue, text width=3.9cm, font=\scriptsize, anchor=west] (fre) at (12.0,-4.0) {
  \textbf{El freno, siempre uno de tres}\\[3pt]
  el RGPD\\ el consentimiento\\ un tercero};
\draw[fl, draw=cMal] (nc.east) -- (fre.west);

% ---------------- la pregunta obligatoria ---------------------------
\node[cajaOk, text width=16.4cm, align=left, font=\scriptsize, anchor=north] (alt) at (6.3,-9.6) {
  {\small\bfseries Y antes de darlo por cerrado: ¿hay una \emph{versión} del conjunto que sí se pueda compartir?}\\[4pt]
  registro con metadatos públicos \quad · \quad versión agregada \quad · \quad
  subconjunto sin las variables sensibles \quad · \quad diccionario, documentación y código};
\draw[fl, draw=cBien] (fre.south) |- (alt.east);

\node[rotulo, text width=15.3cm, align=left, anchor=north west] at (-1.65,-11.4)
  {Tan abierto como sea posible, tan cerrado como sea necesario. La decisión se toma
   \textbf{conjunto por conjunto}: decidir en bloque es el error más frecuente, y casi siempre
   en la dirección restrictiva.};

\end{tikzpicture}
\end{document}
